% Lösungen Einheit 2 von 5 – Tangente als Berührung (zusammenbau v0.1)
\einheitenkopf{Einheit 2 von 5 · Tangente als Berührung}

\erg{13}{a) berührt – Funktionswert und Steigung stimmen überein. \quad b) $f(1) = 0 = g(1)$ und $f'(1) = 3 - 4 = -1$ = Steigung von $g$ – beide Bedingungen erfüllt: $g$ ist Tangente. \quad c) $f(x) = g(x) \Leftrightarrow 2x^2 - 4x + 2 = 2(x - 1)^2 = 0$, Stelle $x = 1$; $f(1) = g(1) = 2$, $f'(1) = g'(1) = 2$; gemeinsame Tangente $y = 2x$ \quad d) $f(2) = 4 = g(2)$ (die Teile stoßen aneinander); $f'(x) = \frac{1}{2}x$, $f'(2) = 1$ = Anstieg von $g$ – gleiche Anstiege, also kein Knick (die Funktionswerte allein zeigen das nicht). \quad e) $f(0) = -16$, $f'(0) = 24$, $t(x) = 24x - 16$; $f(x) - t(x) = x^3 - 9x^2 = x^2(x - 9) = 0$: $x = 0$ ist die Berührstelle, weiterer Punkt bei $x = 9$: $(9 | 200)$ \quad f) Berührpunkt $(2 | 3)$, Tangente $y = 2x - 1$ (ebenso möglich: Berührpunkt $(-2 | 3)$, $y = -2x - 1$) \quad g) An welcher Stelle $u$ die Gerade $y = 3x - 4$ den Graphen von $f$ berührt – dort stimmen Funktionswert und Steigung überein. \quad h) $f(2) = 2$, $f'(x) = (1 - x) \cdot e^{2 - x}$, $f'(2) = -1$; $t(x) = -x + 4$; $t(3) = 1$, $f(3) = 3e^{-1} \approx 1{,}104$; relative Abweichung $\frac{1{,}104 - 1}{1{,}104} \approx 9{,}4\,\%$ (bezogen auf $f(3)$) – ja, weniger als $10\,\%$ \quad i) $t_1(x) = -2x + 4$, $t_2(x) = 2x - 4$; $t_1(2) = 0 = t_2(2)$ – beide gehen durch $S(2 | 0)$. \quad j) Berührbedingung $f'(a) = \frac{10 - f(a)}{5 - a}$ (Steigung der Sichtlinie): $a^2 - 10a - 24 = 0$, $a = -2$ oder $a = 12$; im Bereich liegt $a = -2$: die Sichtlinie berührt den Hang in $(-2 | 3)$, dahinter ist er verdeckt – nicht einsehbar für $-4 \le x < -2$.}

13f)

\begin{ksys}[xmin=-4,xmax=4,ymin=-2,ymax=6,klein] \funktion{0.5*\x^2+1}{f} \gerade{2}{-1}{t} \gerade{-2}{-1}{t_2} \end{ksys}

\erg{14}{a) Tangente in $Q$: $y = f'(u) \cdot (x - u) + f(u)$; $P$ einsetzen: $f(-1) = f'(u) \cdot (-1 - u) + f(u) \Leftrightarrow 3u^4 + 4u^3 - 16u^2 - 32u - 15 = 0 \Leftrightarrow (u + 1)^2 (3u^2 - 2u - 15) = 0$; mit $u > 0$: $u = \frac{1 + \sqrt{46}}{3} \approx 2{,}59$}
\erg{15}{a) Tangente im Punkt $(10 | 30)$ einzeichnen und ihren Schnitt mit der $x$-Achse ablesen: $x = 25$, also nach $25$ Stunden (nicht die Nullstelle des Graphen).}
\erg{16}{a) $g'(4) = 1$, $h'(4) = 2 - 0{,}5 = 1{,}5$ (bei $-4$ symmetrisch $-1$ und $-1{,}5$): $h$ fällt von den Endpunkten aus steiler ab, liegt dort also unterhalb von $g$; wegen $h(0) = 2 > 1 = g(0)$ liegt $h$ in der Mitte oberhalb.}
\erg{17}{a) Jonas prüft nur die Steigung. Es ist $f(2) = 3$, aber $g(2) = 5$ – kein gemeinsamer Punkt, $g$ ist keine Tangente, nur parallel zu ihr. \quad b) Dort ist $f - t$ null, und wegen $f' = t'$ ist auch die Steigung von $f - t$ null: der Graph von $f - t$ berührt die $x$-Achse, die Nullstelle ist doppelt.}
